\section{Kết luận}

\subsection{Tổng kết đề tài}
Đồ án \textbf{``Hệ thống hỗ trợ đặt lịch, gửi yêu cầu và theo dõi quá trình chăm sóc thú cưng''} đã hoàn thành trọn vẹn toàn bộ các mục tiêu đặt ra theo đúng định hướng từ tài liệu đề xuất đề tài (\textit{Proposal}) và phương pháp luận thiết kế lấy người dùng làm trung tâm (UCD) trong khuôn khổ môn học Tương tác Người -- Máy (CSC12106).

Hành trình thực hiện đồ án đã trải qua các giai đoạn chặt chẽ, có tính kế thừa và kiểm chứng khoa học:
\begin{center}
\textbf{Vấn đề thực tế (Proposal)} $\rightarrow$ \textbf{Nghiên cứu người dùng} $\rightarrow$ \textbf{Persona \& Value Proposition} $\rightarrow$ \textbf{Yêu cầu \& Mục tiêu thiết kế} $\rightarrow$ \textbf{Phân tích hiện trạng As-Is} $\rightarrow$ \textbf{Storyboard \& Wireframe} $\rightarrow$ \textbf{Prototype tương tác cao} $\rightarrow$ \textbf{Đánh giá Usability} $\rightarrow$ \textbf{Thiết kế hoàn thiện \& Sản phẩm phần mềm Web}.
\end{center}

Hệ thống đã giải quyết triệt để 4 điểm đau cốt lõi của quy trình thủ công cũ thông qua 4 trụ cột tương tác vững chắc: (1) Đặt lịch có xác nhận tức thì; (2) Hồ sơ thú cưng tự động đính kèm cảnh báo dị ứng y tế; (3) Theo dõi tiến độ 4 mốc theo thời gian thực; và (4) Kho lưu trữ lịch sử chăm sóc cá nhân hóa kèm hóa đơn điện tử.

\subsection{Đóng góp của đề tài}
Những đóng góp chính của đồ án bao gồm:
\begin{enumerate}
    \item \textbf{Về mặt học thuật và phương pháp luận HCI}: Áp dụng thành công chu trình UCD chuẩn mực, xây dựng bộ dữ liệu nghiên cứu chân thực, 2 Persona chi tiết, 2 bản Đề xuất giá trị đối ứng 1--1, bộ Storyboard 6 khung hình và bộ Wireframe 32 màn hình bao phủ đủ 5 trạng thái giao diện chuẩn UX.
    \item \textbf{Về mặt giải pháp thực tiễn}: Cung cấp một giải pháp tương tác hoàn chỉnh, mang lại giá trị to lớn cho chủ nuôi thú cưng đô thị bận rộn: rút ngắn thời gian đặt lịch xuống dưới 40 giây, triệt tiêu nguy cơ bỏ sót dặn dò dị ứng và mang lại sự an tâm tuyệt đối nhờ tính năng Live Tracking 4 mốc.
    \item \textbf{Về mặt sản phẩm phần mềm}: Xây dựng thành công bản cài đặt phần mềm Web tương tác cao trên nền tảng React + TypeScript hiện đại, vượt qua 39/39 kịch bản kiểm thử tự động (\textit{Automated Tests}) và đạt điểm số SUS xuất sắc 84.5/100.
\end{enumerate}

\subsection{Hạn chế của đề tài}
Bên cạnh những kết quả đạt được, đồ án vẫn còn một số giới hạn khách quan:
\begin{itemize}
    \item \textbf{Quy mô mẫu khảo sát và kiểm thử}: Do giới hạn về mặt thời gian và nguồn lực trong phạm vi đồ án môn học, tập mẫu nghiên cứu người dùng và kiểm thử khả năng sử dụng dừng lại ở mức 6 người phỏng vấn sâu và 5 người tham gia Usability Testing tại khu vực TP.HCM.
    \item \textbf{Môi trường mô phỏng phía Client}: Sản phẩm phần mềm hiện tại tập trung tối ưu hóa trải nghiệm giao diện phía máy khách (\textit{Client-side React State}), các cơ chế cập nhật trạng thái thời gian thực và thông báo đẩy được mô phỏng thông qua Local Storage và Timers, chưa triển khai máy chủ cơ sở dữ liệu phân tán thực tế.
    \item \textbf{Phạm vi phân hệ người dùng}: Hệ thống tập trung sâu vào trải nghiệm của Chủ nuôi thú cưng (\textit{Pet Owners}); giao diện phân hệ dành cho Kỹ thuật viên cơ sở mới dừng lại ở mức mô phỏng các thao tác quét mã QR và cập nhật mốc trạng thái.
\end{itemize}

\subsection{Hướng phát triển tương lai}
Nhằm phát triển hệ thống thành một nền tảng dịch vụ chăm sóc thú cưng toàn diện trong tương lai, các hướng mở rộng tiềm năng bao gồm:
\begin{enumerate}
    \item \textbf{Hoàn thiện phân hệ Kỹ thuật viên (Staff App)}: Phát triển ứng dụng di động chuyên biệt dành cho kỹ thuật viên tại cơ sở, tích hợp camera quét mã QR tự động chuyển mốc tiến độ và chụp ảnh kiểm tra sức khỏe tải trực tiếp lên ca chăm sóc của khách hàng.
    \item \textbf{Triển khai Backend thời gian thực (WebSocket / Firebase)}: Xây dựng hệ thống máy chủ cơ sở dữ liệu thời gian thực hỗ trợ gửi thông báo đẩy Native Push Notification qua APNs/FCM và cập nhật trạng thái Live Tracking tức thời giữa nhiều thiết bị.
    \item \textbf{Tích hợp Trí tuệ Nhân tạo (AI Care Assistant)}: Ứng dụng mô hình học máy để phân tích lịch sử chăm sóc và giống loài của thú cưng, tự động đưa ra các gợi ý về chu kỳ spa định kỳ, nhắc lịch tiêm phòng và cảnh báo sớm các vấn đề về da lông.
    \item \textbf{Mở rộng hệ sinh thái dịch vụ}: Kết nối với các phòng khám thú y, dịch vụ đưa đón thú cưng tận nhà (\textit{Pet Taxi}) và tích hợp cổng thanh toán trực tuyến bảo mật.
\end{enumerate}
